\subsection{Setup}%
\label{sub:Setup}

Here are the steps on getting the BC up and running on your local OpenFoam-12.
Open a terminal and do the following:
\begin{enumerate}
  \item Clone the repository to any folder (given that you have access) by typing:\\ \textit{git clone git@gitlab.sintef.no:Hursanay.Fyhn/syntheticturbulenceinlet.git}
  \item Go into the cloned folder by typing\\ \textit{cd syntheticTurbulenceInlet}
  \item Source you local OpenFoam-12 the way you normally source it.
  \item Compile the boundary condition and add it to the list of boundary conditions for your OpenFoam-12 by typing:\\ \textit{wmake}
    \\ (ignore the message about .H files already existing)
  \item You can now use the BC like any other, just remember to source the relevant library by adding this line to your controlDict:\\
    \textit{libs (``libsyntheticTurbulenceInlet.so'');}
  \item Within the cloned directory, you will find a folder titled tutorials. Try out the tutorial cases for the four different geometry modes inside the tutorials folder.
\end{enumerate}

\subsection{Parameters}%
\label{sub:Parameters}

\begin{itemize}
  \item \textit{massFlowMode} (bool) (optional)\\
    Choose if you want to specify a massFlow (with true) or an Umean (with false).\\
    Default: false

  \item \textit{massFlow} (scalar, table, coded or other Function1 types) (required if massFlowMode is true)\\
    Total mass flow rate into the patch.

  \item \textit{Umean} (scalar, table, coded or other Function1 types) (required if massFlowMode is false)\\
    Mean speed into the patch.

  \item \textit{geometryMode} (annulus, disk, rectangle, channel or arbitrary) (required)\\
    Channel is rectangle with PBC in one direction. The purpose of the arbitrary mode is to generate turbulence on a patch with arbitrary shape, do not follow the DNS profiles unlike the other geometry modes due to constant $U$, $R$ and $L$.

  \item \textit{wallLongOrShort} (long or short) (optional)\\
    When \textit{geometryMode} is \textit{channel}, choose if the wall is on the long or the short side.\\
    Default: long

  \item \textit{arbitraryModeRadius} (scalar) (optional)\\
    When \textit{geometryMode} is \textit{arbitrary}, the use can choose the radius within which to generate turbulence.\\
    Default: The largest distance from the patch center.

  \item \textit{arbitraryModeCenter} (vector) (optional)\\
    When \textit{geometryMode} is \textit{arbitrary}, the user can choose the center coordinate.\\
    Default: geometrically averaged center coordinate of the patch.

  \item \textit{Utau} (scalar) (optional)\\
    Friction velocity\\
    Default: Utau is calculated from the expression for a smooth pipe. There is an absolute lower limit that maps the entire DNS profile over the entire width, and it can be forced by e.g. choosing a negative value.

  \item \textit{swirlNr} (scalar) (optional)\\
    Swirl number\\
    Default: 0

  \item \textit{d} (scalar) (optional)\\
    Eddy density (ratio of the sum of eddy volumes to the eddy box volume)\\
    Default: 1

  \item \textit{scale} (scalar) (optional)\\
    A scaling factor for $U'$\\
    Default: 1

  \item \textit{nCellPerEddy} (int) (optional)\\
    Minimum eddy length in units of number of cells\\
    Default: 1

  \item \textit{Lref} (scalar) (optional)\\
    Normalisation factor for integral-length scale\\
    Default: 1
    
  \item \textit{seed} (unsigned int) (optional)\\
    Seed for the random numbers\\
    Default: an arbitrary number

\end{itemize}

\subsection{Example usage}%
\label{sub:Example usage}

\noindent inlet\\
  \{\\
    type            syntheticTurbulenceInlet;\\
    geometryMode    annulus;\\
    massFlowMode    true;\\
    massFlow        2.53e-4;\\
  \}\\

